
\section{Appendix: Extremal length and width}

%For reader's convenience we provide below the necessary background facts concerning the
%extremal length.
%We will use Appendices of \cite{covering lemma,K} for reference to the needed
%background on the conformal length.
%Below we will add a few extra facts that are not contained in the above sources.

Given a family curves $\GG$ on a Riemann surface $S$ and a conformal metric $\mu$ on $S$,
we let $\mu(\gamma)$ be the $\mu$-length of a curve $\gamma\in \Gamma$,
$\mu(\Gamma)$ be the infimum of these lengths,
$\area_\mu$ be the corresponding measure, and
 $\LL(\GG)$ and $\WW(\GG)=\LL(\GG)^{-1}$ be respectively the {\it extremal length} and {\it width} of $\GG$:
see \cite{A} or the Appendices \cite{covering lemma,K} for the precise definitions.
The most basic properties of these conformal invariants, the {\it Parallel} and {\it Series Laws} can
also be found in these sources.


\subsection{Transformation rules}
 Both extremal length and extremal width  are conformal invariants.
 More generally, we have:

\begin{lem}\label{increase}
   Let $f: U\ra V$ be a holomorphic map between two Riemann surfaces,
and let $\GG$ be a family of curves on $U$. Then
$$
      \LL(f(\GG))\geq \LL(\GG).
$$
% Moreover, if $f$ is at most $d-to-1$, then
% $$
%   \LL (f(G)) \leq d\cdot  \LL(\GG).
% $$
\end{lem}

 See Lemma 4.1 of \cite{covering lemma} for a proof.

\begin{cor}\label{two families}
     Under the circumstances of the previous lemma,
let $\HH$ be a family of curves in $V$ satisfying the following lifting
property: any curve $\gamma\in \HH$ contains an arc that lifts to some curve
in $\GG$. Then $\LL(\HH)\geq \LL(\GG)$.
\end{cor}

 See Corollary 10.3 of \cite{K} for a proof.


% \subsection{Extremal distance and Dirichlet integral}

\msk
 Given a compact subset $K\subset \inter U$, the {\it extremal distance}
               $$\LL(U, K)\equiv \mod (U,K)$$
(between $\di U$ and $K$) is defined as $\LL(\GG)$, where $\GG$ is the family of curves
connecting $\di U$ and $K$. In case when $U$ is a topological disk and $K$ is connected,
we obtain the usual modulus $\mod(U\sm K)$ of the annulus $U\sm K$.
%
%\ssk {\it Remark.} $\LL(U,K)$ can also be defined as $\LL(\GG')$ where $\GG'$ is the family of curves
%{\it in} $U\sm K$ connecting $\di U$ to $K$. Indeed, since $\GG\supset \GG'$, $\LL(\GG)\leq \LL(\GG')$.
%Since each curve of $\GG$ overflows some curve of $\GG'$, $\LL(\GG)\geq \LL(\GG')$.
%One can also make a compromise and use the intermediate  family of curves in $U$ connecting $\di U$ to $K$.
%\ssk
%
 We let $\WW(U,K) = \LL^{-1}(U,K)$.

%\begin{lem}\label{modulus transform}
% Let $f: U\ra V$ be a branched covering of degree $N$ between two compact Riemann surfaces with boundary.
%Let $A$ be a compact subset of $\inter U$ and let $B=f(A)$.
%Then
%$$
%   \mod(U,  A)\leq \mod(V, B) \leq N \, \mod(U,A).
%$$
%\end{lem}


%\subsection{More transformation rules}

%The Dirichlet integral formulation allows us to sharpen the lower bound in Lemma \ref{modulus transform}:


\begin{lem}\label{modulus transform-2}
 Let $f: U\ra V$ be a branched covering between two compact Riemann surfaces with boundary.
Let $A$ be an archipelago in $U$,  $B=f(A)$, and
assume that  $f:  A\ra B$ is a branched covering of degree $d$. Then
$$
      \mod(V, B) \geq d\, \mod (U,  A).
$$
\end{lem}

See Lemma 4.3 of \cite{covering lemma} for a proof.


% {\it Remark.} A similar estimate is still valid for an arbitrary compact set $A$,
% and can be proved by approximating $A$ by  archipelagos.

\ssk

%Putting the above two lemmas together
%(or using directly that the Dirichlet integral is transformed as the area
%under branched coverings) we obtain:

%In the case when $f: U\sm A\ra V\sm B$ is a branched covering,
%the modulus is transformed particularly well:

\begin{lem}\label{modulus transform-1}
Let $(U,A)$ and $(V,B)$ be as above,
and let  $f: U\sm A\ra V\sm B$ be a branched covering of degree $N$. Then
$$
      \mod(V, B)= N\, \mod (U,  A).
$$
\end{lem}

See \cite{A} for a proof.



\subsection{Strips and quadrilaterals}

\begin{lem}\label{Pi}
 Let us consider a horizontal strip $\Pi(h)$
and an interval $I=(x, x+a)\subset \R$. We view $\Pi$ as a
quadrilateral with horizontal sides $I$ and $\R+ih$.
Then
$$
      \frac{h}{2a}\leq  \mod \Pi \leq \frac{h}{a},
$$
provided $h/a\leq 1/2$ or $\mod\Pi \leq 1/4$ (for the left-hand side inequality).
\end{lem}

\begin{proof}
  By definition, $\mod\Pi$ is the extremal length of the family of curves connecting
$I$ to  $\R+ih$. This family contains the family $\GG'$ of vertical curves
in the Euclidean rectangle with horizontal sides $I$ and $I+ih$.
Hence $\LL(\GG)\leq \LL(\GG')=h/a$.

To prove the left-hand side inequality, let us consider a Euclidean rectangle
$Q$ with vertices $x-h, x+a+h, x+a+h+ih, x-h+ih$ endowed with the Euclidean metric $\mu$.
Any curve of $\GG$ has $\mu$-length at least $h$. Hence
$$
  \LL(\GG)\geq \frac{h^2}{\area_\mu(Q)}= \frac{t}{1+2t}, \quad where\ t=h/a.
$$
We see that $\LL(\GG)\geq t/2$ for $t\leq 1/2$,
 while $\LL(\GG) >  1/4$ otherwise. The conclusion follows
\end{proof}


\begin{lem}\label{C}
Let $\Pi$ and $I$ be as in the previous lemma.
Let $C= \Pi/\, l\,\Z $
be a cylinder covered by $\Pi$
so that $I$ is embedded into the bottom of $C$.
Then
$$
      \mod \Pi  \geq \frac{1}{2} \min(\mod C,\, 0.5).
$$
% provided $\mod \Pi\leq 1/4$.
\end{lem}

\begin{proof}
Since the covering $\Pi\ra C$ is an embedding on $I$, we have: $a\leq l$.
 Then by the previous lemma we obtain:
$$
    \mod \Pi \geq \frac{h}{2a} \geq \frac {h}{2l}= \frac{1}{2} \mod C.
$$
\end{proof}



\comm{********
Similar estimates hold for a cylinder in place of the strip:

\begin{lem}\label{C}
 Let us consider a flat cylinder $C$ of height $h$.
Let $B$ and $T$ be the ``bottom'' and the ``top'' circles of $C$.
Let  $I\subset B$ be an interval of length $a$.
Then
$$
      \frac{h}{2a}\leq  \ed(I, T) \leq \frac{h}{a},
$$
provided $h/a\leq 1/2$ (for the left-hand side inequality).
\end{lem}

It allows us to estimate the modulus of a holomorphic rectangle
inside an annulus:

\begin{prop}\label{quadr in cyl}
  Let $Q$ be a quadrilateral with base $I$ and top $J$,
and let $C$ be a flat cylinder as in the previous lemma.
Let $\phi: Q\ra A$ be a holomorphic map
which continuously extends to $\bar Q\sm \di J$.
Assume  that $\hat \phi$ embeds $I$ into $B$,
maps the vertical sides of $Q$ to $B$, and maps $\inter J$ to $T$.
Then
$$
   \frac{1}{2} \mod A \leq \mod Q \leq \mod A,
$$
provided ....
\end{prop}

\begin{proof}
    Let  $\Pi=\{z: 0< \Im z< h\}  $
be a strip realization of the universal covering of $C$
such that $\R$ covers the bottom of $C$.
Then $\phi$ lifts to a holomorphic map $\hat\phi: Q\ra \Pi$
that embeds $I$ to $\R$, maps the vertical sides of $Q$ to $\R$,
and maps  $J$ to $\R+ih$.

Let $\HH$ be the family of vertical curves in $Q$
that do not pass through $\di J$,
and let $\GG$ be the family of curves in $\Pi$ that connect $I$ to $\R+ih$.
Then any curve of $\HH$ maps to a curve of $\GG$, and vice versa,
any curve of $\GG$ lifts to a vertical curve in $Q$.
Thus $\hat \phi(\HH)=\GG$.
By Lemma \ref{increase},
$$\mod(Q)= \LL(\HH) \leq \LL(\GG).$$
\end{proof}
********************}

\subsection{Holomorphic and embedded annuli}

Let $S$ be a hyperbolic Riemann surface with boundary with a preferred component $\si$ of $\di S$.
We assume that $S$ has finite topological type and is not the punctured disk.
A {\it holomorphic annulus}  in $\AAA$ is a holomorphic map $A: \A(1,r) \ra S$
that extends to a homeomorphism  $\phi: \T\ra \si$.
We let $\mod A = \mod \A(1,r)$.
% Let $\AAA$ stand for the family of holomorhic annuli in $S$.


The family of holomorphic annuli contains a subfamily of {\it embedded annuli}.
Among embedded annuli,  there is an  annulus $A_*$ of maximal modulus,
which has nice special properties. Namely, let us uniformize $A_*$ by a flat cylinder
$C= \Pi(h)/ \Z$. Then the quadratic differential $dz^2$ on $C$ is the pull-back of some
quadratic differential $q$ on $S$. Moreover, the uniformization $C\ra A_*$ extends continuously
to the upper boundary $C^+= \R+ih /\Z$ of $C$ (minus finitely many points corresponding to the punctures
of $S$),
and induces there an equivalence relation
$\tau_k : \alpha_k \ra \alpha_k'$, where $(\tau_k)$ is a finite family of  isometries between pairs of
disjoint arcs in $C^+$. The images of these arcs, $\la_k = i(\alpha_k)=i(\alpha_k')$,
are {\it horizontal separatricies} of $q$.
(It is a version of Strebel's Theorem, see e.g., \cite[\S 11]{GL}).

\begin{lem}\label{emb vs hol}
  For any holomorphic annulus $A: \A(1,r)\ra S$, we have:
$$
   \mod A\leq 16 \mod A_*.
$$
\end{lem}

\begin{proof}
  Let us consider a family $\GG$ of non-trivial proper curves $\gamma$ in $S$ that begin in  $\si$.%
\footnote{``Non-trivial'' means that $\gamma$ cannot be pulled to $\si$ through a continuous
family of proper curves.} Then  any curve $\gamma\in \GG$ contains an initial segment $\gamma'$
that lifts to a vertical curve in $\A(1,r)$. By Corollary \ref{two families},
\begin{equation}\label{upper bound}
  \mod A\leq \LL(\GG).
\end{equation}

Let us now take any conformal metric $\mu$ on $S$, and let $l= \mu(\GG)$.
For any vertical curve $\de$ in $A_*$, two possibilities can occur:

\ssk\nin $\bullet$ $\de$ ends on $\di S$. Then $\de \in \GG$ and hence $\mu(\de)\geq l$.

\ssk\nin $\bullet$  $\de$ ends on some separatrix $i(\alpha_k)$.
Then there is another vertical curve $\la$
in $A_*$ that ends at the same point as $\de$.
The concatenation of $\de$ and $\la$ is a curve of family $\GG$.
Hence one of the curves, $\de$ or $\la$, is ``long'',
i.e., it has $\mu$-length at least $l/2$.
\ssk

It follows that at least one half of the vertical curves in $A_*$ are long.
Let $I\subset C^+$ be the set of endpoints of $i^{-1}$(long curves).
We can now proceed as in the classical Gr\"otztsch estimate.
By the Cauchy-Schwarz Inequality,
$$
 h\, \area_\mu(C) = \area(C) \int_C \area_\mu\, dx\, dy   \geq  \left( \int_I dx \int_{0}^h \mu(x,y) dy\right)^2
  \geq \left( \frac{l}{4} \right)^2,
$$
which implies
$$
  \LL_\mu(\GG)=l^2/\area_\mu(C) \leq 16 h = 16 \mod A_*.
$$
Since this is valid for any conformal metric $\mu$, we conclude that $\LL(\GG)\leq 16\mod A_*$.
Together with (\ref{upper bound}), this gives us the desired estimate.
\end{proof}
